% ============================================================================
% 4. BASELINES — Journal version (expanded with implementation details)
% ============================================================================
\section{Baseline Implementations}
\label{sec:baselines}

We implement six baselines spanning the spectrum from static text injection
to interactive tool-augmented reasoning.  All baselines are zero-shot and
use identical question prompts; only the context injection method varies.

\subsection{Flat Text (FT)}

The entire schema is serialized as a text document listing all tables,
columns, data types, foreign key relationships, and lineage edges.  This
context is prepended to the question, simulating a developer who has
loaded the complete DDL documentation into a long-context LLM.  No
retrieval or filtering is applied.

\subsection{Vector-RAG (VR)}

Schema elements (tables, edges, documentation) are embedded using
Sentence-BERT~\cite{reimers2019sentencebert} and indexed in a FAISS vector
store~\cite{johnson2019faiss}.  At query time, the $k=15$ most relevant chunks are retrieved and
injected as context.  This tests whether dense retrieval can identify the
relevant structural fragments for topology questions.

\subsection{Graph-Augmented (GA)}

The question is first parsed to identify mentioned table names.  A 3-hop
BFS neighborhood is extracted from the schema graph centered on these
anchor tables, and the subgraph (with FK and lineage edges) is serialized
as structured context.  This provides the LLM with local topology
information without requiring full schema memorization.  For global
questions lacking explicit anchor tables (e.g., ``How many connected
components?''), the full graph is provided as context.

\subsection{Tool-Use (TU)}

The LLM is given access to purpose-built graph query tools:
\texttt{get\_neighbors(table, edge\_type)},
\texttt{shortest\_path(src, dst)},
\texttt{get\_connected\_component(table)},
\texttt{get\_lineage\_chain(table, direction)}, and
\texttt{list\_tables()}.  The LLM can make up to 3 sequential tool calls
per question, interleaving reasoning with graph algorithm execution.
This tests whether interactive graph access can overcome the limitations
of static context injection.

\subsection{ReAct-Code (RC)}

The LLM generates and executes Python code using NetworkX directly on the
schema graph object, with up to 5 code execution rounds per question.
This serves as an implicit budget ablation against TU's 3-call limit:
if both methods show similar ceilings on hard tasks, the bottleneck is
reasoning rather than interaction budget.  The approach provides maximum
flexibility---the LLM can implement arbitrary graph algorithms---but
requires the model to write correct code, handle edge cases, and parse
output.

\subsection{Oracle}

The gold-standard algorithmic output is injected directly into the prompt.
For a \texttt{join\_path} question, the Oracle context includes the exact
shortest path; for \texttt{combined\_impact}, it includes the complete set
of affected tables.  Oracle scores establish the upper bound achievable by
each model (97\% EM; the 3\% gap is due to list-ordering canonicalization
and alternative valid path formats), isolating language understanding from
graph retrieval.

\subsection{Tier-2 Baselines}

For Tier-2 value-level questions, we implement three adapted baselines:
\textbf{FT$_{v2}$} (schema DDL context only),
\textbf{GA$_{v2}$} (schema + lineage summaries + value samples), and
\textbf{TU$_{v2}$} (interactive lineage/data query tools including
row-level access).
